\chapter{테니스 치는 뉴런}
% full version: chapters/22_dishbrain.tex

칩 위의 살아 있는 뉴런으로 한 실험 가운데 가장 유명한 것은 “접시 속의 뇌” DishBrain이다. 전극이 2만 6000개 달린 칩 위에서 기른 뉴런 약 100만 개가 가장 단순한 컴퓨터 테니스를 쳤다. 공이 화면을 날아다니고, 패들은 위아래로 움직인다.

\section*{게임은 어떻게 만들어져 있나}

공은 전극 8개를 통해 네트워크로 들어간다. 그 가운데 어느 전극이 딸깍하느냐가 공의 높이를 알린다. 얼마나 자주 딸깍하느냐가 공이 얼마나 먼지를 알린다. 공이 먼 벽에 있으면 초당 4번이고, 패들로 다가올수록 점점 잦아져 40번까지 올라간다. 출력은 전극 영역 두 곳이다. 한 영역의 딸깍은 패들을 위로, 다른 영역의 딸깍은 아래로 움직인다. 컴퓨터는 10밀리초마다 딸깍을 세어 패들을 움직인다. 이 10밀리초는 배양이 아니라 시험대 컴퓨터의 클록이다. 네트워크 자체는 모든 신경망처럼 비동기로 일한다.

\pencilfig{22}{\pencilart{22}}{공과 패들, 그리고 패들을 움직이는 뉴런 100만 개.}

\section*{교사}

도파민은 전혀 없다. 히트가 나면 네트워크는 짧고 완전히 예측 가능한 신호를 받는다. 전극 8개 전부에 한꺼번에 0.1초 동안이다. 미스가 나면 무작위 자리에 무작위 순간으로 4초 동안 어지러운 자극을 받고, 잠시 쉰 뒤 새 서브가 들어온다.

\section*{결과}

20분짜리 세션 동안 살아 있는 배양의 랠리는 길어졌고 놓친 서브는 줄었다. 가장 잘된 것은 피드백이 온전할 때였다. 히트와 미스 뒤의 신호 대신 그냥 조용하기만 했을 때도 나아지긴 했지만 덜했다. 피드백이 없으면 나아지지 않았다. 사람 뉴런이 생쥐 뉴런보다 조금 잘 쳤다. 반면 며칠에 걸친 학습은 믿을 만하지 않은 것으로 드러났다.

\section*{네트워크는 왜 배우는가}

저자들은 결과를 이렇게 설명한다. 네트워크는 자기 감각을 예측 가능하게 만들려고 한다. 미스는 혼돈을, 히트는 질서를 가져오므로, 네트워크는 질서가 늘어나도록 바뀐다. 그러나 데이터와 똑같이 잘 맞는 더 현실적인 설명도 있다. “미스 뒤에는 흔들고, 히트 뒤에는 건드리지 마라.” 흔들 때마다 연결이 무작위로 조금씩 뒤섞인다. 네트워크가 자주 미스를 내는 설정은 계속 뒤섞이고, 잘되는 설정은 손대지 않은 채 남는다. 네트워크는 미스가 적은 곳에 저절로 머무르게 된다. 아무도 가르치지 않았다. 그저 더는 흔들지 않을 뿐이다. 우리의 단순한 모델에서 이 규칙은 실제로 히트의 비율을 끌어올린다.

\pencilfig{130}{%
% scrambler: miss probability p over one setting of the network (valley in the middle); a miss kicks the setting
% at random, a hit leaves it alone, so the time spent at a setting is proportional to 1/p (6x more in the valley here)
\draw[pen] (0,1.2) -- (8.2,1.2);
\draw[penlite=pcRed] plot[smooth] coordinates {(0.0,2.46) (0.8,2.46) (1.6,2.46) (2.0,2.44) (2.4,2.41) (2.8,2.32) (3.2,2.15) (3.6,1.90) (4.0,1.62) (4.4,1.44) (4.8,1.44) (5.2,1.62) (5.6,1.90) (6.0,2.15) (6.4,2.32) (6.8,2.41) (7.2,2.44) (7.6,2.46) (8.0,2.46)};
\node[lbl, text=pcRed, anchor=south] at (7.55,2.55) {미스 많음};
\node[lbl, text=pcRed, anchor=west] at (5.3,1.45) {적음};
% kicked balls on the slopes
\draw[dotpen=pcBlue, wash=pcBlue] (1.3,2.68) circle (0.12);
\draw[arr=pcGray, densely dashed] (1.35,2.82) to[bend left=50] (2.35,2.72);
\draw[arr=pcGray, densely dashed] (1.22,2.82) to[bend right=50] (0.4,2.72);
\draw[dotpen=pcBlue, wash=pcBlue] (6.3,2.45) circle (0.12);
\draw[arr=pcGray, densely dashed] (6.25,2.59) to[bend right=50] (5.45,2.25);
\node[lbl, anchor=south] at (1.3,3.05) {미스 뒤 흔들기};
% the ball left alone in the valley
\draw[dotpen=pcBlue, wash=pcBlue] (4.6,1.66) circle (0.12);
\node[lbl, anchor=south] at (4.6,1.95) {흔들지 않음};
% where the network spends its time (proportional to 1/p)
\fill[pcGreen, opacity=0.22] (0,0) -- plot[smooth] coordinates {(0.0,0.18) (0.8,0.18) (1.6,0.18) (2.0,0.19) (2.4,0.19) (2.8,0.21) (3.2,0.24) (3.6,0.33) (4.0,0.55) (4.4,0.98) (4.8,0.98) (5.2,0.55) (5.6,0.33) (6.0,0.24) (6.4,0.21) (6.8,0.19) (7.2,0.19) (7.6,0.18) (8.0,0.18)} -- (8.0,0) -- cycle;
\draw[penlite=pcGreen] plot[smooth] coordinates {(0.0,0.18) (0.8,0.18) (1.6,0.18) (2.0,0.19) (2.4,0.19) (2.8,0.21) (3.2,0.24) (3.6,0.33) (4.0,0.55) (4.4,0.98) (4.8,0.98) (5.2,0.55) (5.6,0.33) (6.0,0.24) (6.4,0.21) (6.8,0.19) (7.2,0.19) (7.6,0.18) (8.0,0.18)};
\draw[pen] (0,0) -- (8.2,0);
\node[lbl, text=pcGreen!60!black, anchor=west] at (5.3,0.6) {네트워크가 머무는 곳};
\node[lbl, anchor=north] at (4.1,-0.03) {네트워크 설정};
}{미스는 네트워크 설정을 무작위로 흔들고, 히트는 그대로 둔다. 미스가 많은 곳에는 네트워크가 머물지 않는다. 미스가 적은 곳에서는 흔들림이 멈추고 네트워크는 거기 남는다. 아무도 거기로 이끌지 않았는데도.}

두 설명을 어떻게 가려낼까? 아직 아무도 하지 않은 실험이 필요하다. 미스 뒤에 어지러운 신호가 아니라, 길이와 세기가 같은 \emph{예측 가능한} 신호를 넣는 것이다. 그래도 네트워크가 배운다면, 핵심은 예측 가능성이 아니라 그저 “히트 뒤”와 “미스 뒤”의 차이다. 이 실험은 아직 수행할 사람을 기다리고 있다.

이제 우리 여정이 끝났다. 20와트짜리 기계는 단순한 부품으로 조립되어 있지만, 그 부품을 정확히 어떻게 엮어 생각을 만드는지 우리는 아직 일부만 이해한다. 어쩌면 다음 장은 여러분이 쓰게 될지도 모른다.

\fullversion{22}
